\documentclass[a4paper,11pt]{article}

\usepackage[utf8]{inputenc}
\usepackage[T1]{fontenc}
\usepackage[italian]{babel}
\usepackage{amsmath,amssymb}
\usepackage[margin=2.5cm]{geometry}
\usepackage{enumitem}
\usepackage{titlesec}

% --- checklist: casella vuota (da fare) o spuntata (confermato) ---
\newlist{checklist}{itemize}{2}
\setlist[checklist]{label=$\square$, leftmargin=1.5em, itemsep=5pt, topsep=4pt}
\setlist[itemize,2]{label=$-$, leftmargin=2.2em, itemsep=1pt, topsep=1pt}
\newcommand{\done}{\item[$\blacksquare$]}

\titleformat{\section}{\normalfont\Large\bfseries}{\thesection}{0.6em}{}
\titleformat{\subsection}{\normalfont\large\bfseries}{\thesubsection}{0.6em}{}

\title{Argomenti di Geometria e Algebra Lineare\\da recuperare per i corsi del II anno}
\date{\today}
\author{Liam Ferretti}

\begin{document}
\maketitle
\thispagestyle{empty}


% ============================================================
\section{Analisi Vettoriale --- confermato (prof.ssa Francesca De Marchis)}
% ============================================================
\begin{checklist}
  \item Prodotto matrice/matrice e matrice/vettore
  \item Prodotto scalare tra vettori \textit{(verrà comunque reintrodotto a lezione)}
  \item Calcolo del determinante di una matrice
  \item Concetto di rango di una matrice
  \item Matrici (o forme quadratiche) definite positive, negative, semidefinite
        positive/negative e indefinite, con caratterizzazione in termini di
        autovalori \textit{(verranno comunque richiamate a lezione)}
  \item Calcolo del Jacobiano
  \item Ripasso Analisi in $\mathbf{R}^n$
\end{checklist}

\vspace{0.4cm}

% ============================================================
\section{Meccanica Analitica --- confermato (prof.ssor Mauro Papinutto)}
% ============================================================
\begin{checklist}
  \item Condizioni di indipendenza lineare di un insieme di vettori
  \item Applicazioni lineari e loro nucleo
  \item Rango di una matrice
  \item Matrici simmetriche
  \item Matrici ortogonali
  \item Problema agli autovalori e autovettori
  \item Diagonalizzazione di matrici simmetriche
  \item Forme quadratiche e segnatura
\end{checklist}

\vspace{0.4cm}
\noindent\textit{Nota: da Termodinamica e Laboratorio e da Laboratorio di Fisica
Computazionale~1 non erano emersi collegamenti diretti significativi.}

\end{document}